% Copyright 2026 The dhnt Authors. Licensed under Apache-2.0.
% Build: pdflatex dhnt && pdflatex dhnt   (run twice for references)
\documentclass[11pt]{article}

\usepackage[margin=1in]{geometry}
\usepackage{amsmath,amssymb}
\usepackage{booktabs}
\usepackage{enumitem}
\usepackage{xcolor}
\usepackage{listings}
\usepackage[hidelinks]{hyperref}

\lstset{basicstyle=\ttfamily\small,breaklines=true,frame=single,
  columns=fullflexible,keepspaces=true,xleftmargin=4pt,framexleftmargin=4pt}

\newcommand{\dhnt}{\textsc{dhnt}}
\newtheorem{proposition}{Proposition}
\newtheorem{definition}{Definition}

\title{\dhnt{}: Verifiable, Self-Healing Agent Skills\\
via a Contract-Carrying Controlled Language}

\author{Qiang Li\\
\texttt{liqiang@gmail.com}\\[2pt]
\small (affiliation to be added)}

\date{\today\\[4pt]\small Preprint. Reference implementation:
\url{https://github.com/dhnt/dhnt}}

\begin{document}
\maketitle

\begin{abstract}
Agent ``skills'' --- the reusable, model-followed procedures now standardized by
the \texttt{SKILL.md} format and adopted across dozens of agentic tools --- are
today \emph{natural-language instructions a model interprets}. Their behavior
varies across models and across runs, their success is judged by the model
itself (so a wrong-but-plausible result passes silently), and they have no
stable identity, so a fix discovered on one run is re-discovered on the next.
We present \dhnt{}, a controlled language and runtime in which a skill is a
\emph{contract-carrying, content-addressed, typed artifact}. A \dhnt{} skill has
a canonical form (with a content hash as its identity), a machine-checkable
\emph{success contract}, a typed and bounded \emph{effect} declaration, and
portable run \emph{attestations}. We argue and demonstrate that the contract
functions as a \emph{deterministic oracle}: it makes success verification
independent of the executing model, bounds and audits side effects, and---most
distinctively---enables \emph{sound runtime self-improvement}, in which a
contract-verified repair is cached and folded back into the skill as an
environment-guarded branch, then reused across agents without re-incurring the
failure. \dhnt{} round-trips to the \texttt{SKILL.md} standard, so generic tools
consume a \dhnt{} skill unchanged while \dhnt{}-aware runtimes additionally
execute and verify it. We describe the language, its determinism and soundness
properties, a Go reference implementation, and preliminary results; and we
outline a controlled evaluation of cross-model convergence, effect containment,
and self-healing amortization.
\end{abstract}

\section{Introduction}

The agentic-tooling ecosystem has rapidly converged on a shared unit of reusable
capability: the \emph{skill}. The \texttt{SKILL.md} format---YAML frontmatter, a
Markdown body, optional bundled scripts, and progressive disclosure---was
published as an open standard in late 2025 and adopted by dozens of tools within
months. A skill, in this model, is \emph{prose}: instructions a language model
reads and interprets.

Prose is a remarkably effective interface, but it has four structural weaknesses
as a unit of \emph{reliable} capability:
\begin{enumerate}[topsep=2pt,itemsep=1pt]
\item \textbf{Non-determinism across models.} The same skill yields different
behavior on different models/providers.
\item \textbf{Non-determinism across runs.} Sampling makes even one model vary
run to run.
\item \textbf{Unverified success.} ``Did it work?'' is judged by the model; a
plausible-but-wrong result passes silently.
\item \textbf{No identity, no memory.} Two phrasings of one procedure are
unrelated strings; a fix found on one run (e.g.\ adapting a script to a platform
quirk) is re-discovered on the next, because there is nowhere to durably store a
\emph{trustworthy} fix.
\end{enumerate}

We argue these are not deficiencies of skills \emph{per se} but of skills
expressed only as prose, and that the missing ingredient is a \emph{deterministic
contract}: a machine-checkable statement of what success means, decided
identically regardless of who (or which model) executed the skill. With such a
contract, success becomes verifiable rather than judged; side effects can be
bounded and audited; and---the key consequence---a learned fix can be cached and
reused \emph{soundly}, because the contract certifies it.

\paragraph{Contributions.}
\begin{enumerate}[topsep=2pt,itemsep=1pt]
\item \textbf{\dhnt{}, a controlled language for skills} in which a skill is a
typed AST over a closed multilingual glossary, with a canonical, content-addressed
form and bidirectional projection to human languages (\S\ref{sec:lang}).
\item \textbf{A contract-and-effect type discipline}: success is a conjunction of
typed predicates; side effects are a bounded lattice; both are enforced at run
verification (\S\ref{sec:lang},~\S\ref{sec:attest}).
\item \textbf{A precise determinism account}: we separate determinism of the
artifact, of verification, and of execution, and state in what sense weak and
strong executors \emph{converge} (\S\ref{sec:det}).
\item \textbf{Portable, re-checkable attestations} of skill runs
(\S\ref{sec:attest}).
\item \textbf{Sound runtime self-improvement}: contract-gated, context-keyed
repair caching and branch accretion, with a soundness argument
(\S\ref{sec:heal}).
\item \textbf{Zero-cost interoperability} with the \texttt{SKILL.md} standard via
graceful degradation (\S\ref{sec:interop}), and a Go reference implementation
(\S\ref{sec:impl}).
\end{enumerate}

\section{Background and Related Work}\label{sec:related}

\paragraph{Agent skills and protocols.} The \texttt{SKILL.md} standard
\cite{agentskills} defines the prevailing skill artifact. Tool- and
agent-interoperability standards---the Model Context Protocol \cite{mcp},
Agent-to-Agent \cite{a2a}, and \texttt{AGENTS.md}, now under the Linux
Foundation's Agentic AI Foundation---standardize \emph{transport and discovery},
not \emph{procedure semantics}. \dhnt{} is complementary: it specifies what a
skill \emph{is} and how to verify it, and is designed as an optional profile of
\texttt{SKILL.md}.

\paragraph{Skill acquisition.} ReAct \cite{react}, Toolformer \cite{toolformer},
and Voyager \cite{voyager} let agents plan with and accumulate skills. These
skills are unverified and re-discovered; \dhnt{} contributes verified,
content-addressed skills with sound reuse.

\paragraph{Contracts and neurosymbolic verification.} Design by Contract
\cite{meyer} inspires recent work placing a contract layer around LLM calls
\cite{dbcagent}, verifying instruction-following neuro-symbolically
\cite{nsverify}, synthesizing specifications \cite{specsynth}, and generating
verified policies \cite{verafi}. \dhnt{} extends this from a per-call wrapper to
a full language, interlingua, and runtime with content identity and self-healing.

\paragraph{Controlled natural language.} Attempto Controlled English
\cite{ace} translates a restricted English into first-order logic for executable
specifications; Gherkin/BDD pairs readable scenarios with executable steps.
\dhnt{}'s controlled language adds a content-addressed canonical form,
multilingual projection, effect typing, and attestation.

\paragraph{Capabilities and attestation.} Effect/permission manifests and
pre-action authorization for agents \cite{permmanifest} motivate bounding agent
side effects. Verifiable-build attestation (in-toto \cite{intoto}, SLSA) provides
the closest analogue to \dhnt{}'s run receipts---signed, re-checkable records of
\emph{how} an artifact was produced---here applied to skill execution.

\section{Overview}\label{sec:overview}

Consider a release skill: ``ensure the tests pass and the tag is signed.'' As
prose, a model interprets and executes it; whether it succeeded is the model's
opinion, and if a step's script breaks on a particular shell version, the model
may fix it---then forget the fix by the next invocation.

In \dhnt{}, the same skill is the artifact in Figure~\ref{fig:skill}. Its
\emph{contract} (\texttt{tests-green}, \texttt{tag-signed}) is checked
mechanically; its \emph{effect cap} (\texttt{\{read, write\}}) bounds what any
execution may do; its canonical form hashes to a stable \emph{identity}; and a
run emits an \emph{attestation} anyone can re-check. When a step fails in a given
environment, a model proposes a fix that is accepted only if it still satisfies
the original contract within the original effect cap; the fix is then folded into
the skill as an environment-guarded branch and cached, so the next run---by any
agent sharing the store---reuses it.

\begin{figure}[t]
\begin{lstlisting}
skill release
  effects read write              # blast radius (P3)
  ensure tests-green              # success contract (P1)
  ensure tag-signed
  step announce print value text
  when dirty:                     # flow control / playbook
    step stop abort
\end{lstlisting}
\caption{A \dhnt{} skill (English projection). The same artifact has a canonical
$a$--$z$ form that hashes to its identity and re-parses losslessly, and projects
into any registered language.}
\label{fig:skill}
\end{figure}

\section{The \dhnt{} Language}\label{sec:lang}

\dhnt{} is layered. Free prose (L0) is normalized by a model into a controlled,
glossary-locked form (L1), which deterministically encodes to a canonical
$a$--$z$ string (L1.5), which parses to a typed AST (L2), which an executor
realizes against bindings (L3). Validity is \emph{transpilability}: an expression
is valid iff it round-trips through L1.5.

\paragraph{Identity (content addressing).} The canonical form is a pure function
of the AST; its hash is the skill's identity. Equal canonical form $\Rightarrow$
equal identity, on any machine. This yields cross-model equality, deduplication,
caching, and lineage---``content addressing for procedures'' \cite{merkle}.

\paragraph{Contract.} A skill carries a set of \emph{checks}, each a typed
predicate over the glossary. A run is \emph{valid} only if every check evaluates
to true (the boolean is native to the language's numeral system). The contract
generalizes an opaque verify-command into a portable, typed predicate.

\paragraph{Effects.} Each primitive/predicate declares effects drawn from a
closed lattice $\{\textsf{read}, \textsf{write}, \textsf{net}, \textsf{spend},
\textsf{destroy}, \textsf{time}\}$; a skill declares an upper bound
(\emph{effect cap}). A run is valid only if its observed effects are a subset of
the cap.

\paragraph{Steps, latitude, branches, composition.} Steps are leaf actions; each
carries a \emph{latitude} (exact vs.\ judgement) that tunes how deterministic it
is. A step may instead be a \emph{branch} (\texttt{when}/\texttt{else}),
turning a linear runbook into a branching playbook. A step's primitive may
resolve to another skill (composition); a content-addressed library audits the
call graph and enforces that a caller's effect cap contains each callee's.

\paragraph{Bidirectionality.} The AST projects into any registered language and
parses back, so a human may edit any language face and have it re-normalize to
the single canonical identity.

\section{Determinism}\label{sec:det}

\dhnt{} does not make a language model deterministic. It makes three other things
deterministic and pushes model non-determinism into explicit, gated places.

\begin{enumerate}[topsep=2pt,itemsep=1pt]
\item \textbf{Artifact.} Canonical form, AST, identity, and projections are pure
functions: the same skill yields the same bytes and identity everywhere.
\item \textbf{Verification.} The contract, the effect-cap check, and the
attestation are deterministic functions of a run's observed results; whether a
run satisfied the spec is decided identically for every executor.
\item \textbf{Execution.} A step bound to a code leaf is deterministic; a step
bound to a model is as non-deterministic as that model---but is always gated by
verification.
\end{enumerate}

\begin{proposition}[Outcome convergence]
Let $s$ be a skill with contract $C$ and effect cap $E$, executed by tiers
$t_1,t_2$ producing observed results $r_1,r_2$ and effect sets $F_1,F_2$. The
validity verdict $V(s,r,F) \equiv C(r)\wedge (F\subseteq E)$ is a deterministic
predicate; hence $t_1$ and $t_2$ \emph{converge on the success criterion} ($V$
agrees) whenever they reach states satisfying $C$ within $E$, even if their
execution paths and outputs differ.
\end{proposition}

This is the precise sense of ``leveling the playing field'': a weak executor with
a strong contract and a strong executor with the same contract agree on
\emph{success}, though not necessarily on output. The choice of how much of a
skill is code (deterministic) vs.\ model (gated) is the latitude dial.

\paragraph{Caveat.} ``Deterministic'' is relative to \emph{(skill, world-state,
bindings)}, as for any program; and the L0$\to$L2 normalizer is a model call,
hence non-deterministic \emph{at authoring time}---but its output is validated and
then frozen as a deterministic artifact with an identity, as when a model writes
code that is thereafter run deterministically.

\section{Verifiable Attestation}\label{sec:attest}

A run emits an \emph{attestation}: the skill identity, the executor tier, the
predicates that passed/failed, the observed effects, an optional world-diff, and
a validity bit computed (not asserted) as $C(r)\wedge(F\subseteq E)$. Anyone
holding the skill and the attestation can re-check consistency. Attestations make
``consistently close'' auditable: success is provable, not hoped for, analogous
to build attestation \cite{intoto}.

\section{Sound Runtime Self-Improvement}\label{sec:heal}

The contract enables what prose skills cannot do safely: \emph{cache a fix}.

\paragraph{Loop.} Given a skill, a context fingerprint (declared environment
probes), a store, and a model:
(1) if a verified variant is cached for \emph{(identity, context)}, run it and
re-verify; (2) else run the skill as written; (3) else ask the model for a
repair. A repair is \emph{accepted only if its steps satisfy the original
contract within the original effect cap}---the verification skill grafts the
original contract and cap onto the candidate's steps, so the model may change the
\emph{how} but never the \emph{what}. An accepted repair is cached, and
\emph{folded} into a new skill version as an environment-guarded branch
(\texttt{when} env-matches \texttt{then} fix \texttt{else} original). Folds nest,
so the skill accretes one arm per environment.

\begin{proposition}[Promotion soundness]
Every variant promoted to the cache satisfies the original contract within the
original effect cap in the context under which it was promoted; therefore reusing
a cached variant (after re-verification) never accepts a result the original
specification would reject, and never exceeds the original blast radius.
\end{proposition}

This is exactly the property a prose skill lacks: it has no identity to key a
cache on and no oracle to certify a cached fix, so it must re-interpret each run.
Where the folded version lives is a deliberate policy: a host-local overlay
(self-healing host) by default; promotion to a shared catalog is a separate,
human-reviewed step (the change is model-authored and may be context-specific).

\section{Interoperability with the Agent Skills Standard}\label{sec:interop}

\dhnt{} renders a skill as a standard \texttt{SKILL.md}: frontmatter
(\texttt{name}, \texttt{description}, \texttt{executor: cnl}) plus a
self-contained body (constraints from the effect cap, steps---branches rendered
as if/else, success criteria from the contract) and the embedded canonical form.
A generic tool reads the description and follows the prose; a \dhnt{}-aware
runtime executes the canonical form and attests. A tool that does not know
\dhnt{} ignores the unknown executor. Thus one artifact degrades gracefully, and
adoption requires \emph{no change} to existing tools (\S\ref{sec:eval}).

\section{Implementation}\label{sec:impl}

The reference implementation is a dependency-light Go module: a core
\texttt{skills} package (canonical form, parser/lineariser, glossary, contract,
effects, attestation, composition library, normalizer, exporter, adaptation
runtime), a PTY executor leaf for driving terminal tools, and a \texttt{dhnt}
CLI (\texttt{export}, \texttt{run}, \texttt{normalise}, \texttt{promote}). The
core has no LLM dependency; a model is injected as a \texttt{Completer}, with an
adapter that uses an existing agent CLI as the model. The spec is published with
the implementation.

\section{Evaluation}\label{sec:eval}

We report preliminary, reproducible results from the reference implementation and
outline the controlled study required for full claims. \emph{(Numbers marked
\textsc{[tbd]} are pending the controlled runs of \S\ref{sec:eval}.)}

\paragraph{Illustrative results (implemented).}
\begin{itemize}[topsep=2pt,itemsep=1pt]
\item \textbf{Interoperability.} A \dhnt{}-generated \texttt{SKILL.md} drives
real terminal tools over a pseudo-terminal: an interactive coding CLI end-to-end
(spawn, handle a confirmation prompt via a branch, query, exit), and five agent
CLIs via a token-free version probe, each contract-verified.
\item \textbf{Self-healing amortization (synthetic).} In a controlled harness, a
skill that fails in each new environment is repaired once and thereafter reused:
across five invocations spanning two environments, the system performs exactly
two repairs, the rest served from the learned, accreted version.
\item \textbf{Determinism.} Round-trip and idempotence of the canonical form, and
re-checkability of attestations, hold by construction and are covered by tests.
\item \textbf{Safety.} Repairs that exceed the declared effect cap are rejected
and not cached.
\end{itemize}

\paragraph{Reproducible harness.} A suite (\texttt{eval/}) implements E1--E5; its
hermetic configuration produces illustrative figures---e.g.\ effect containment
with \emph{zero} false-accepts across the candidate set; self-healing with two
repairs over five invocations spanning two environments; a deterministic tier
with a single distinct verdict across repeats while a noisy-but-correct tier
varies in output but not verdict. Publication numbers require running the same
harness with live providers and a real task suite, below.

\paragraph{Planned controlled study.}
\begin{description}[topsep=2pt,itemsep=1pt,leftmargin=2.4em]
\item[E1 Cross-model convergence.] One skill, $N$ models (including small ones);
measure success rate and output variance \emph{with vs.\ without} the contract
gate. Hypothesis: the validity verdict's cross-model variance collapses while raw
output variance does not. \textsc{[tbd]}
\item[E2 Determinism/repeatability.] Variance across repeated runs for code-leaf
vs.\ model-leaf skills. \textsc{[tbd]}
\item[E3 Self-healing on a real task suite.] Cross-platform shell skills;
report repairs vs.\ invocations and amortized cost. \textsc{[tbd]}
\item[E4 Effect containment.] Injected unsafe repairs; report containment recall
/ false-positive rate. \textsc{[tbd]}
\item[E5 Interop breadth.] Generated \texttt{SKILL.md} executed across $\geq 3$
standard tools unchanged. \textsc{[tbd]}
\end{description}
Baselines: plain prose skills; a Design-by-Contract per-call wrapper
\cite{dbcagent}. Metrics: success rate, verdict variance, repair amortization,
containment recall, interop pass rate.

\section{Limitations}\label{sec:limits}

The glossary is closed (mitigated by a deterministic loan-word rule for names);
rich, multi-branch playbooks authored directly from long prose are not yet
robust; full-screen interactive TUIs require a capability-answering terminal;
effect enforcement currently \emph{detects and attests} over-cap effects rather
than \emph{preventing} them before they land (a sandboxed leaf is future work);
and the controlled evaluation (\S\ref{sec:eval}) is in progress.

\section{Conclusion}

A deterministic contract is the missing oracle that turns agent skills from
interpreted prose into verifiable, content-addressed, self-healing, portable
artifacts---and it can be added \emph{inside} the standard the ecosystem already
adopted, at near-zero integration cost. \dhnt{} is a controlled language and
runtime that realizes this, and a candidate verifiable profile of the Agent
Skills standard.

\begin{thebibliography}{99}
\small
\bibitem{agentskills} Anthropic. \emph{Agent Skills} (\texttt{SKILL.md}) open
standard, 2025. \url{https://www.anthropic.com/news/skills}
\bibitem{mcp} Anthropic. \emph{Model Context Protocol}, 2024.
\url{https://modelcontextprotocol.io}
\bibitem{a2a} Google. \emph{Agent-to-Agent (A2A) Protocol}, 2025 (Linux
Foundation Agentic AI Foundation).
\bibitem{react} S. Yao et al. \emph{ReAct: Synergizing Reasoning and Acting in
Language Models}. arXiv:2210.03629, 2022.
\bibitem{toolformer} T. Schick et al. \emph{Toolformer}. arXiv:2302.04761, 2023.
\bibitem{voyager} G. Wang et al. \emph{Voyager: An Open-Ended Embodied Agent with
Large Language Models}. arXiv:2305.16291, 2023.
\bibitem{meyer} B. Meyer. \emph{Applying ``Design by Contract''}. IEEE Computer,
25(10), 1992.
\bibitem{dbcagent} \emph{A DbC-Inspired Neurosymbolic Layer for Trustworthy Agent
Design}. arXiv:2508.03665, 2025.
\bibitem{nsverify} \emph{Neuro-Symbolic Verification on Instruction Following of
LLMs}. arXiv:2601.17789, 2026.
\bibitem{specsynth} \emph{Seeking Specifications: The Case for Neuro-Symbolic
Specification Synthesis}. arXiv:2504.21061, 2025.
\bibitem{verafi} \emph{VERAFI: Verified Agentic Financial Intelligence through
Neurosymbolic Policy Generation}. arXiv:2512.14744, 2025.
\bibitem{permmanifest} \emph{Permission Manifests for Web Agents}.
arXiv:2601.02371, 2026.
\bibitem{ace} N. E. Fuchs, R. Schwitter. \emph{Attempto Controlled English
(ACE)}. CLAW, 1996.
\bibitem{intoto} S. Torres-Arias et al. \emph{in-toto: Providing Farm-to-Table
Guarantees for Bits and Bytes}. USENIX Security, 2019.
\bibitem{merkle} R. C. Merkle. \emph{A Digital Signature Based on a Conventional
Encryption Function}. CRYPTO, 1987.
\end{thebibliography}

\end{document}
